\chapter{Database engines as a systems comparison}
\label{evod42:chapter}

An artifact-promotion record has been appended to a log, but the process crashes
before the serving metadata page is updated. Which version is active after
restart? Now move the crash earlier, before the log reaches durable storage.
The answer should change. Merely saying ``use transactions'' does not specify
either case.

Database engines provide concrete precedents for separating intent, legal
physical plans, execution, buffers, effects and recovery. We will borrow those
mechanisms with explicit premises. A token generator is not a relational query,
and a local dictionary simulating crash cuts is not a production database.
The chapter retains the supplied comparison and WAL example while making
semantic eligibility, record binding and the crash model precise.

\section{Logical requirements delimit physical alternatives}
\label{evod42:plans}
\index{logical and physical plans}\index{cost model!hard constraints}
Let $G$ denote a required logical observation contract and
$\mathcal P(G)$ the physical implementations that preserve it. A cost optimizer
may choose
\[
 p^\star=\arg\min_{p\in\mathcal P(G),\ M_p\le M_{max},\ R_p\le R_{max}} J(p).
\]
The legal set is not a penalty term. An implementation with a different
tokenizer, artifact, numerical policy or state-continuation rule does not
become equivalent because it is cheap. Chapter 41's verifier and lowerer
determine legality before placement and scheduling optimize predicted cost.

Our fixture offers two abstract native-v1 alternatives and a cheaper
attention-v1 alternative. These are stipulated contracts, not measurements of
DOGMA versus Hermon quality. With latency in milliseconds, payload residency
in MiB, risk a dimensionless policy proxy, and weights
$w=(1,0.1,10)$, the illustrative score is
\[
 J=L_{ms}+0.1M_{MiB}+10R.
\]
Weights carry reciprocal units into one score. Risk remains subject to a hard
ceiling in addition to its score term. It is not a calibrated failure
probability in this example.

\Plate{03-plans}{The cheapest candidate is ineligible. Contract and hard resource filters precede ranking, so the fixture chooses native-B.}{evod42:plansplate}
\begin{center}\input{\Generated/alternatives.tex}\end{center}
\Listing{choose}
Selection validates finite nonnegative metrics and weights, unique names,
resource ceilings and risk range. It breaks ties by name and fails explicitly
when no legal alternative remains. A NaN must not silently win or evade a
comparison. A model substitution that changes quality is a new task-level
choice, not automatically an equivalent physical-plan rewrite.

\section{Volcano and Cascades transfer structure, not equivalence}
\label{evod42:precedents}
\index{Volcano}\index{Cascades}\index{prefix-dependent generation}
Volcano's operator interface separates execution mechanics from the semantics
of individual item operations. Its exchange operator encapsulates parallel
data movement~\cite{evod42:volcano}. The transferable lesson is to keep module
boundaries stable while changing execution mechanism. A native state step,
attention decode, retrieval read and symbolic operation may share lifecycle
entry points without sharing an internal tensor representation.

Cascades organizes alternative expressions and physical requirements through
memoized optimization tasks and rule-guided search~\cite{evod42:cascades}.
Evolutor could similarly search legal implementations under required properties,
but needs an equivalence relation for each rule. A memo entry keyed only by
prompt text could incorrectly merge different owners, model versions or
random streams. We implement no Cascades memo or search algorithm here.

An autoregressive token step is prefix-stateful. Its selected output can become
the next input, changing future state and work. It is not an iterator over
independent rows. Nor is a database operator necessarily stateless: joins and
aggregates retain state. The useful distinction concerns the declared algebra
and legal transformations, not the simplistic claim that databases have no
state while models do.

In a serving runtime, batching compatible requests can resemble vectorized
operator execution, but each lane still needs its own local position, owner,
artifact and completion fence. Chapter 34 proved such a lift for a restricted
recurrence. The database analogy does not prove attention batching or make
Chapter 36's unfinished physical allocation interface a passing dependency.

\section{Buffer management needs object-specific validity}
\label{evod42:buffers}
\index{buffer manager}\index{pin!physical lifetime}\index{index!authority}
A buffer manager tracks residency, ownership, pins and replacement. These
mechanisms are relevant to weights, native recurrence state, attention KV,
durable records and retrieved context. Their reconstruction rules differ.
An immutable weight tensor can be reloaded from a verified artifact; a native
state needs a complete checkpoint or exact replay; Hermon KV needs the exact
ordered prefix, model and learned positions; a forgotten event may no longer
be eligible for restoration at all.

\Plate{02-objects}{Representation and recovery semantics remain attached to each object type. Allocating equal-size pages does not make their contents interchangeable.}{evod42:objects}

Pins protect physical lifetime, not logical validity. A pinned buffer with the
wrong artifact is still wrong; a logically valid handle cannot be reclaimed
while device work reads it. Chapter 32 implements generational native-pool
ownership and delayed reclaim, while Chapter 35 specifies ordered layered KV.
When Chapter 36 settles, its page-sharing and eviction tests must preserve
both sets of relevant lifetime constraints. This chapter does not substitute
a generic byte-page demo for those implementations.

Database index ideas also transfer with limits. Exact key/range access and
approximate vector retrieval support different query contracts. An index is
an access path, not the authoritative evidence object. Chapter 40 checks current
tombstones and all summary ancestors before returning a stale index hit.
Promoting the hit into context does not give it policy authority.

\section{Transactions belong to particular durable mutations}
\label{evod42:transactions}
\index{transactions!external effects}\index{materialization}
Atomicity, consistency, isolation and durability can describe metadata writes,
memory updates and an artifact-pointer switch when backed by a real
transaction system. They do not by themselves describe sampling. A generated
token may be computed privately, accepted into a request's output stream, or
delivered to a client; those are separate events with different rollback
possibilities.

Isolation should likewise name the shared resource. Concurrent durable updates
to a tenant's memory might need serializable transactions. Separate request-local
KV tensors need ownership and lifetime isolation. A phrase such as ``serializable
generation'' does not specify either contract. External actions may support
an idempotency key or compensation, but a database rollback cannot reverse an
already sent message or completed physical action.

Materialization is a useful comparison for immutable prefixes, compiled
artifacts and derived summaries. Reuse requires exact semantic identity,
tenant/policy scope, source availability and the relevant artifact versions.
An engine-native prefix cache and a semantic document cache are different
materialized objects with different keys and invalidation rules.

\section{Derive a small write-ahead recovery contract}
\label{evod42:wal}
\index{write-ahead logging}\index{log sequence number}\index{page LSN}
ARIES is a primary precedent for WAL, log sequence numbers and recovery that
repeats history with rollback machinery~\cite{evod42:aries}. Our reference is
\emph{committed-redo-only}: every assignment becomes an independently committed
unit when its log prefix is flushed. There are no concurrent transactions,
uncommitted page writes, undo, compensation log records, checkpoints, torn
pages, media loss or real \texttt{fsync}. It is not an ARIES implementation.

The crash model has a volatile append-only log, durable prefix boundary $D$,
durable scalar pages $data[k]$, and page sequence numbers $P_k$. Appending an
assignment creates immutable record $(\ell,k,v)$ with the next LSN $\ell$.
Flushing advances $D$ only over existing contiguous records. A page write is
legal only when
\begin{equation}
 1\le\ell\le D,\qquad log[\ell]=(\ell,k,v),\qquad P_k\le\ell.
 \label{evod42:walrule}
\end{equation}
The first condition is log-before-page ordering. The second binds the proposed
page contents to the record, preventing a caller from using an unrelated
durable LSN as permission for arbitrary data. The third prevents a stale page
write from overwriting a newer assignment.

\Listing{WalStore.write_page}
The package's simpler check tested only $\ell\le D$. It could therefore
accept a different value or key under a valid record number. The strengthened
negative controls reject both, reject flushes beyond log end and reject stale
per-key writes. A hash or LSN is meaningful only when it binds the actual
state transition being certified.

\Plate{01-wal}{The durable log frontier precedes the durable page frontier. A crash in between requires redo; a crash before log flush loses the volatile assignment.}{evod42:walplate}

On crash, discard the volatile tail above $D$ and keep durable pages.
Recovery scans retained records in LSN order and writes a record only when
its page LSN is older.
\Listing{WalStore.recover}

\begin{proposition}[Redo correctness and idempotence]
Assume initially empty pages, only legal page writes under
Equation~\ref{evod42:walrule}, and a retained correct contiguous log prefix
$1,\ldots,D$. Recovery leaves each key equal to its latest assignment in that
prefix, or absent if it has no assignment. Running recovery again changes
neither data nor page LSNs.
\end{proposition}
\begin{proof}
For key $k$, let $\ell_k$ be its maximum retained LSN. Legal page writes
ensure any existing $P_k$ names an actual retained record for $k$ and
$P_k\le\ell_k$. Scanning in increasing order either skips an already newer
page or advances it to the record being scanned. By the time $\ell_k$ is
visited, the page equals that assignment and its LSN is $\ell_k$. A second
scan finds every record's LSN at most the current page LSN and skips it.
Keys absent from the retained log cannot have a legal page write.
\end{proof}
Corruption, missing log ranges or forged page LSNs violate the premises.
Production storage needs checksums, validated recovery metadata and a real
durability mechanism. The theorem does not turn an in-memory copy into a
disk-crash guarantee.

\section{Crash cuts need an independent expected state}
\label{evod42:controls}
\index{crash-cut testing}\index{redo!idempotence}
The second recovery oracle simply applies last assignment per key over the
durable prefix. It never invokes the page-LSN algorithm.
\Listing{replay_prefix}

\begin{center}\input{\Generated/crashes.tex}\end{center}
Cuts zero through six occur before work and after each append, flush and page
write for v1 and v2. After append without flush, recovery cannot see the new
assignment. After flush without page write, recovery supplies it. After page
write, redo is a no-op for that record. The generated table and assigned
expected answer vector test these distinctions rather than only asserting
that recovery returns \emph{something}.

Randomized schedules append assignments to three keys, flush varying prefixes,
write some eligible pages and crash after every iteration. Recovered maps
must equal the independent prefix specification, and second recovery must be
idempotent. These tests exercise a stipulated fault model; they do not measure
power-loss behavior, torn writes or an operating system's durability semantics.

\section{Recovery, replay and stale statistics are different problems}
\label{evod42:limits}
\index{recovery versus replay}\index{optimizer!statistics}
Recovery restores a consistent durable metadata state. Replay reproduces
declared execution observations under captured inputs, artifacts, RNG and
numerical policy. A recovered serving pointer can be correct without
regenerating every historical token. Conversely, numerically replaying a
request does not prove its external tool effect was committed only once.

Cost statistics can be wrong while a plan remains legal. Relevant serving
statistics include prompt/output lengths, native-state bytes, prefix reuse,
speculative acceptance, batch occupancy, backend queues and retrieval quality.
They need workload identity, collection time and uncertainty. Misestimation
can cause a poor memory/latency choice; it must not expand semantic eligibility.
The figures here contain assigned arithmetic, not measured workload statistics.

Four limits keep the comparison bounded: prefix-dependent token generation
does not inherit relational rewrite laws; sampling is not ACID commit; native
and KV objects have model-specific validity; and task quality/risk needs its
own evidence beyond a physical cost estimate. Database engines are useful
because they make such contracts explicit, not because using their vocabulary
makes an inference runtime a database.

\section{Twelve worked exercises}
\label{evod42:exercises}
\begin{enumerate}
\item Why reject a cost 0.11 candidate when another costs 5.0?
\textbf{Answer:} The cheaper candidate has the wrong logical contract and is
outside $\mathcal P(G)$.
\item Compute native-A's score for $(5\,ms,4\,MiB,.03)$.
\textbf{Answer:} $5+.1(4)+10(.03)=5.7$ under the declared weights.
\item What transfers from Volcano? \textbf{Answer:} Stable operator boundaries
and encapsulated execution/parallelism. A shared interface does not equate
native recurrence state and layered attention KV.
\item Is every database operator stateless? \textbf{Answer:} No; joins and
aggregates retain state. The legal algebra and observations determine rewrites.
\item Is an unpinned but logically valid page reclaimable?
\textbf{Answer:} Only after all relevant owners, references and in-flight
consumers release it under the actual allocator contract.
\item What three checks authorize the miniature page write?
\textbf{Answer:} Its LSN exists and is durable, its key/value matches that
record, and it does not regress the page's LSN.
\item Crash after append but before flush: what survives?
\textbf{Answer:} The new assignment is outside the durable prefix and is lost.
Earlier durable assignments remain recoverable.
\item Crash after flush but before page write: what happens?
\textbf{Answer:} Redo installs the retained assignment if the page LSN is older.
\item Why is redo idempotent? \textbf{Answer:} After first recovery each
page's LSN is its latest durable assignment; all older/equal records are skipped.
\item Does this reference implement transaction rollback?
\textbf{Answer:} No. Every flushed assignment is committed; undo and multi-record
atomicity are excluded. It cannot model an uncommitted transaction's dirty page.
\item Can model rollback reverse an external action?
\textbf{Answer:} No. A tool needs its own idempotency, reconciliation or
compensation protocol; an irreversible effect remains irreversible.
\item Design a real durability experiment. \textbf{Rubric:} Define storage
and flush guarantees, inject crashes at named cuts, reopen actual files,
validate log/page integrity and committed-prefix behavior, repeat recovery,
and distinguish application-process failure from power loss. Dictionary copying
alone does not qualify.
\end{enumerate}

\section{Vocabulary and the governance boundary}
\Term{Physical alternative}: an implementation satisfying a fixed logical
contract. \Term{WAL}: durable-log-before-durable-page ordering under a specified
fault model. \Term{Recovery}: restoring consistent durable state, distinct from
replaying model execution. Chapter 43 applies these lessons to immutable
candidate lineage, controlled promotion and validated rollback.
